\input{header.tex}
\input{newcommands.tex}


\title{Notensätze mit MusiX\TeX{} erstellen}
\author{Fortbildungsmaterial für MINT-Lehrkräfte}
\date{13. Mai 2025}
\pagestyle{plain}
\usepackage{musixtex} % nur in Musik-Dokumenten: definiert \cup als Note um

\begin{document}

\maketitle

MusiX\TeX{} ist ein hochentwickeltes Paket für professionellen Notensatz in \LaTeX{}. Es ermöglicht Musiklehrkräften, rechtlich unbedenklich eigene Noten zu setzen, ohne auf urheberrechtlich geschützte Vorlagen angewiesen zu sein. Der Name MusiX\TeX{} steht für \glqq{}Music eXtended \TeX{}\grqq{} und unterstreicht die Erweiterung der \TeX{}-Funktionalität um musikalische Notation. Diese Handreichung führt Sie systematisch in die Grundlagen des Notensatzes mit MusiX\TeX{} ein und zeigt, wie Sie komplette Musikstücke für Ihren Unterricht erstellen können.

%%%%%  KAPITEL 1  %%%%
\section{Grundlagen und Paketeinbindung}

Um MusiX\TeX{} zu verwenden, müssen Sie zunächst die notwendigen Pakete einbinden und die Grundstruktur verstehen:

\begin{lstlisting}[style=latex]
\usepackage{musixtex}
\end{lstlisting}

Die zentrale Umgebung für Musiknotation ist die \texttt{music}-Umgebung, in der alle musikalischen Elemente definiert werden:

\begin{lstlisting}[style=latex]
\begin{music}
    % Alle Notationen hier
\end{music}
\end{lstlisting}

Jedes MusiX\TeX{}-Dokument folgt einem strukturierten Aufbau. Zunächst werden die Instrumente und deren Eigenschaften definiert, dann folgt der eigentliche Notensatz zwischen \texttt{\textbackslash startextract} und \texttt{\textbackslash endextract}:

\begin{lstlisting}[style=latex]
\begin{music}
	\startextract
		... Hier stehen die Noten
	\endextract
	\startpiece
		... Hier stehen die Noten
	\endpiece
\end{music}
\end{lstlisting}

\begin{lstlisting}[style=latex]
\begin{music}
\startextract
\Notes\wh{c}\en                        % Ganze Note c
\Notes\hu{d}\hu{e}\en                  % Zwei halbe Noten d und e
\Notes\qu{f}\qu{g}\qu{h}\qu{}\en      % Vier Viertelnoten
\Notes\qu{g}\qu{f}\hu{e}\en            % Gemischte Notenwerte
\endextract
\end{music}
\end{lstlisting}

\begin{ausgabe}
\begin{music}
\startextract
\Notes\wh{c}\en                        % Ganze Note c
\Notes\hu{d}\hu{e}\en                  % Zwei halbe Noten d und e
\Notes\qu{f}\qu{g}\qu{h}\qu{}\en      % Vier Viertelnoten
\Notes\qu{g}\qu{f}\hu{e}\en            % Gemischte Notenwerte
\endextract
\end{music}
\end{ausgabe}

Die wichtigsten Parameter sind \texttt{\textbackslash instrumentnumber}, das die Anzahl der Instrumente festlegt, \texttt{\textbackslash setname}, das den Namen eines Instruments definiert, \texttt{\textbackslash setstaffs}, das die Anzahl der Notensysteme pro Instrument bestimmt, und \texttt{\textbackslash generalmeter}, das die Taktart für alle Instrumente setzt.

%%%%%  KAPITEL 2  %%%%
\section{Grundlegende Notenwerte und einfache Melodien}

MusiX\TeX{} verwendet spezielle Befehle für verschiedene Notenwerte. Die Grundnotenwerte werden mit prägnanten Abkürzungen dargestellt:

\begin{lstlisting}[style=latex]
\begin{music}
\parindent10mm
\instrumentnumber{1}
\setstaffs1{1}
\generalmeter{\meterfrac44}
\startextract
\Notes\wh{c}\en                 			% Ganze Note c (wh=whole note)
\Notes\hu{d}\en                 			% Halbe Note d nach oben (hu=half note up)
\Notes\hl{e}\en                  			% Halbe Note e nach unten (hl=half note low)
\Notes\qu{f}\en      				% Viertelnote f, Hals nach oben (qu=quarter note up)
\Notes\ql{g}\en      				% Viertelnote g, Hals nach unten (ql=quarter note low)
\Notes\qa{a}\en							% Viertelnote a, Hals automatisch (qa=quarter note automatic)
\bar
\Notes\cu{c}\cu{d}\cu{e}\cu{f}\cu{g}\cu{a}\cu{b}\cu{c}\en  % Acht Achtelnoten
\endextract
\end{music}
\end{lstlisting}

\begin{ausgabe}
\begin{music}
\parindent10mm
\instrumentnumber{1}
\setstaffs1{1}
\generalmeter{\meterfrac44}
\startextract
\Notes\wh{c}\en                      % Ganze Note c (wh=whole note)
\Notes\hu{d}\en                  % Halbe Note d nach oben (hu=half note up)
\Notes\hl{e}\en                  % Halbe Note e nach unten (hl=half note low)
\Notes\qu{f}\en      % Viertelnote f, Hals nach oben (qu=quarter note up)
\Notes\ql{g}\en      % Viertelnote g, Hals nach unten (ql=quarter note low)
\Notes\qa{a}\en      % Viertelnote a, Hals automatisch (qa=quarter note automatic)
\bar
\Notes\cu{c}\cu{d}\cu{e}\cu{f}\cu{g}\cu{a}\cu{b}\cu{c}\en  % Acht Achtelnoten
\endextract
\end{music}
\end{ausgabe}

Die Notenwerte folgen einem logischen System: \texttt{\textbackslash wh} steht für \textit{whole} (ganze Note), \texttt{\textbackslash hu} für \textit{half} (halbe Note), \texttt{\textbackslash qu} für \textit{quarter} (Viertelnote) und \texttt{\textbackslash cu} für \textit{croche} (französisch für Achtelnote). Der Befehl \texttt{\textbackslash en} beendet eine Notengruppe, während \texttt{\textbackslash bar} einen Taktstrich setzt.

Oktaven werden durch verschiedene Notationsformen dargestellt. Die Grundoktave verwendet Kleinbuchstaben, tiefere Oktaven Großbuchstaben und höhere Oktaven Apostrophe:

\begin{lstlisting}[style=latex]
\begin{music}
\parindent10mm
\instrumentnumber{1}
\setstaffs1{1}
\generalmeter{\meterfrac44}
\startextract
\Notes\qu{C}\qu{c}\qu{'c}\qu{''c}\en  % Verschiedene Oktaven
\bar
\Notes\qu{D}\qu{d}\qu{'d}\qu{''d}\en  % D in verschiedenen Oktaven
\endextract
\end{music}
\end{lstlisting}

\begin{ausgabe}
\begin{center}
\begin{music}
\parindent10mm
\instrumentnumber{1}
\setstaffs1{1}
\generalmeter{\meterfrac44}
\startextract
\Notes\qu{C}\qu{c}\qu{'c}\qu{''c}\en  % Verschiedene Oktaven
\bar
\Notes\qu{D}\qu{d}\qu{'d}\qu{''d}\en  % D in verschiedenen Oktaven
\endextract
\end{music}
\end{center}
\end{ausgabe}
\end{document}